\section{Contribution, limitations and further questions}

The mixed-product theorems identify attained maxima when there are zero, one or two inconsistent centered singletons. Their proofs keep the exact depleted moment, merge colors through positive partition weights, and handle arbitrary multiplicities through affine indicator identities. For one distinguished inconsistent singleton, an exact kernel difference exchanges a missing conflicting neighbor for one extra prescribed match. This also improves the general inconsistent-singleton bound. The scalar consequences recover the spectral estimate with a smaller singleton factor and an optimal universal prefactor. The diagonal results determine the original cubic threshold, the parity transition, a hierarchy of further corrections and sharp first uniform errors for both the square-root and full-range approximations.

The base recurrence, its normalized parity form and the leading coefficients belong to Collins and Nagatsu \cite{CN}. The mixed estimate itself is already part of Bordenave's spectral-gap method \cite{Bordenave}. The special-function identities and the saddle method are classical. The attained envelope at fixed support and singleton counts is the principal finite extremal result. The diagonal theorems give explicit uniform consequences and elementary proofs in the probabilistic normalization, including the necessity and bounded-comparability characterization of the cubic scale. Classical uniform asymptotic theory supplies another route to the corrections. We do not claim a new special-function method or base historical novelty on a mismatch with only one prior parameter regime.

Bordenave and Zhu \cite[Lemma 2.1]{BZ} retain the original singleton factor $3k/\sqrt N$ and explicitly bound the raw factorial ratio by $\sqrt e$ in their equation (2.3). That elementary bound is a known ingredient here. Our exact refinement retains $\mu_h(n,a-h)$ in addition to $F_n(a-h)$. The factor $\Gamma_{n,a,h}$ is a scalar upper bound on that moment and is smaller than $\beta$ when $1\le h<a$. Theorems~\ref{thm:mixed} and \ref{thm:oneconflict} give exact maxima for $b=0$ and $b=1$, respectively. The latter theorem also replaces the general factor $2^bK_{h,a-h}$ by $2^{b-1}K_{h,a-h-1}$. Theorem~\ref{thm:twoconflict} further subtracts a nonnegative kernel term and attains the resulting maximum at $b=2$. No exact maximum is claimed for $b>2$. A same-row clique need not be extremal for larger $b$. At $n=a=b=8$ and $h=0$, use four disjoint rows, each requesting its own two columns, and center all eight cells. Expanding by rows gives
\[
n^8\E P=\sum_{j=0}^4\binom4j(-2)^j\frac{8^j}{(8)_j}
=\frac{1121}{105}>7,
\]
whereas the eight-cell same-row clique has scaled absolute expectation seven.

The depleted polynomial itself is also a classical specialization. If $d_h^{(r)}(x)$ denotes the higher-order derangement polynomial defined by $(1-t)^{-r}e^{xt}$ in \cite[equation (2.22)]{KKJK}, then
\[
\mathcal P_h(n,s)=(-1)^h d_h^{(n-s-h+1)}(-n).
\]
Indeed, the indicator expansion gives $\mathcal P_h(n,s)=\E(n-V)^h$ for a gamma variable $V$ of shape $n-s-h+1$ and unit scale. Its moment generating function proves the identity. The order is positive on $h+s\le n$ and depends on $h$. Thus neither the polynomial family nor this elementary connection is asserted to be new. The extremal conclusions concern the full class of mixed cell patterns.

Related centered-matrix and permanent inequalities already bound the same diagonal moments. For example, set $B_n=nI-J$, where $J$ is the all-ones matrix. For the matrix elementary symmetric polynomial of Han and Niles-Weed \cite[equation (2), Theorem 1.2]{HNcentered}, summing over injections gives
\[
e_h(B_n)=\binom nh^2\mu_h(n,0),\qquad
\|B_n\|_F=n\sqrt{n-1}.
\]
Their theorem therefore implies
\[
0\le\mu_h(n,0)\le
\frac{(8e^2)^h(n-1)^{h/2}}{\binom nh}\qquad(1\le h\le n).
\]
The permanent expansion in their earlier work \cite[equation (29), Lemma 5.5 in the inspected preprint]{HNmixtures} also contains these moments: the coefficient of $t^h$ in
$n^n\operatorname{Perm}((J+tB_n)/n)/n!$ is $\binom nh\mu_h(n,0)$.
These are genuine overlaps. The inspected statements do not give the exact maximum over our mixed class, whose proof additionally uses positive depletion weights, color merging and attainment.

The Laguerre representation makes the parameter regime precise: the degree is $k$, the parameter is $\alpha=-N-1$, and the argument is $-N$. Thus $-\alpha/k\to\infty$ throughout the sublinear range. Kuijlaars and McLaughlin \cite[Theorem 7.1 and Remark 5.4]{KM}, and the regime announced by Dai and Wong \cite{DW}, treat a finite limiting ratio greater than one. Their parameter ranges overlap the interior proportional-degree part of our full-range formula. Theorem 1 of Dai, Ismail and Wang \cite{DIW} assumes $\alpha>0$ when $\alpha/k\to\infty$. In the inspected preprint of Dunster, Gil and Segura \cite[Section 1, equation (1.6) and Cases 1a, 1b and 2]{DGS}, $\alpha=(k+1/2)(a^2-1)$ with $a\ge0$, excluding $\alpha=-N-1$ for $k\le N$. These comparisons concern those specific statements. They do not exclude other classical uniform expansions.

In particular, the Whittaker connection formulas put the problem at positive argument $N$, with parameters $\kappa=N/2-k$ and $\mu=(N+1)/2$. For every fixed $c<1$, the range $k\le cN$ lies in the uniform domain of Olver's parabolic-cylinder approximations \cite{Olver,DLMFWhittaker}. The odd moment involves a difference of the two solutions, so cancellation must be retained when transporting individual relative errors. There is no parameter-domain obstruction to this route in the sublinear range.

The underlying integral connection can also be checked without special-function notation. Expanding $(N-s)^k$ and integrating each monomial gives the exact identity
\begin{equation}\label{eq:gammaidentity}
Q_k=\frac1{(N-k)!}\int_0^\infty
e^{-s}s^{N-k}(N-s)^k\,ds.
\end{equation}
Indeed, its expansion is
$\sum_{j=0}^k\binom kj(-1)^{k-j}N^j(N-j)!/(N-k)!$,
which is the indicator formula for $Q_k$. Splitting the integral at $N$ and putting $s=N(1\mp t)$ yields
\begin{align*}
Q_k&=\frac{N^{N+1}e^{-N}}{(N-k)!}
\bigl(J_++(-1)^kJ_-\bigr),\\
J_+&=\int_0^1 e^{Nt}t^k(1-t)^{N-k}\,dt,\qquad
J_-=\int_0^\infty e^{-Nt}t^k(1+t)^{N-k}\,dt.
\end{align*}
For $k<N$, both positive integrals have their stationary point at $t=\sqrt{k/N}$. Thus even orders add two contributions and odd orders subtract them. Standard Laplace expansion of this classical representation provides another route to the action and first correction. The positive partition proof used above makes the low-order patch and the explicit error constants transparent.

At $k=N$, the identity with a truncated exponential is classical, as recorded in \cite[Section 1]{DMO}. The contribution of Theorem~\ref{thm:global} is a self-contained formula that patches fixed orders, the parity transition and the endpoint while keeping $(N)_k$ exact. Its optimal global error is $C_*N^{-1/3}+o(N^{-1/3})$. Appendix~\ref{sec:globalsharpproof} gives the full parity profile, controls both tails and proves uniqueness of the limiting maximizing scale. The sharp $N^{-1/4}$ rate in Section~\ref{sec:rates} concerns the different, quartically corrected hyperbolic approximation. These explicit probabilistic refinements are not claimed to be historically independent of classical asymptotic theory.

A bounded source search found no exact duplicate of the attained mixed maximum or the complete collection of uniform probabilistic statements. This is a search outcome, not a priority certificate. Some journal bodies remained inaccessible, as recorded in the accompanying source audit. The elementary integral calculation above is not presented as a newly discovered special-function identity.

The inequalities are proved symbolically for their stated parameter ranges. Exact finite computations check normalization, the depleted recurrence and the conditioning examples. Numerical plots illustrate the asymptotic regimes. Neither kind of computation is used as a substitute for a general proof.

The remaining graph-level task is to combine the sharper cell-pattern estimate with circuit enumeration and tangle remainders while retaining the degree dependence. A better estimate for one fixed product does not automatically improve a spectral theorem: the patterns that dominate the full sum must be identified first. The present paper does not assert a new second-eigenvalue bound or an end-to-end audit of either random-graph proof.
